\chapter{Démonstrations des résultats d'estimations géométriques} \label{chap:demo_estimations_geo}

Cette annexe est consacrée aux démonstrations des résultats d'estimations géométriques énoncés dans les \Cref{prop:estimations_tailles_plat_vs_courbe,prop:mesure_sigmacourbe}. En suivant \cite[Sec. 4.2]{dziuk_finite_2013}, nous considérons une approximation affine par morceaux d'une surface lisse $\surflisse \in \Cont^2$. Nous donnons des estimations de l'erreur géométrique produite par l'approximation de la surface lisse $\surflisse$ par une surface polyédrique $\surfpoly$. Un simplexe courbe $\primal{K} = \projvariete \primalplat{K}$ est paramétré sur un simplexe plan, comme illustré dans la \Cref{fig:estimation_projection}.

\begin{figure}[H]
    \centering
    \tikzset{
    ddfv/.style={
        scale=1.5,
        every path/.style={thick,line join=round, cap=round},
        ptsq/.style 2 args={
          rectangle,
          fill=##1,
          inner sep=2pt,
          label={[text=##1]##2}
        },
        ptsqbord/.style 2 args={
          rectangle,
          draw=##1,
          inner sep=2pt,
          label={[text=##1]##2}
        },
        ptcl/.style 2 args={
          circle,
          fill=##1,
          inner sep=1.5pt,
          label={[text=##1]##2}
        },
        ptclbord/.style 2 args={
          circle,
          draw=##1,
          inner sep=1.5pt,
          label={[text=##1]##2}
        },
        mailleprimale/.style={
            myblue,
            fill=mylightblue
        },
        mailledualeint/.style={
            myred,
            fill=mylightred
        },
        mailledualebord/.style={
            myred,
            fill=mylightred!50!white
        },
        maillediamantint/.style={
            mygreen,
            fill=mylightgreen
        },
        maillediamantbord/.style={
            mygreen,
            fill=mylightgreen!50!white
        }
    }
}

\begin{tikzpicture}[
  ddfv,
  scale=6
]

\def\largeur{3}
\def\xL{1}

\def\f(#1){0.2*sin(pi*#1 r)}
\def\fp(#1){0.2*pi*cos(pi*#1 r)}

% \def\f(#1){0.2*sin(pi*#1^3 r)}
% \def\fp(#1){0.2*3*pi*#1^2*cos(pi*#1^3 r)}

\draw[myblue] (0.5,0.2) -- (0.5,{\f(\xL)+0.2}) node[midway, above] {$\projvariete \primal{K}$};
\draw[myblue] (0,0) -- (\xL,{\f(\xL)}) node[midway, below] {$\primal{K}$};

\draw plot[domain=-0.1:{\xL+0.1},samples=20,smooth] ({\x},{\f(\x)}) node[above right] {$\variete{M}$};

\node[ptsq={myblue}{}] at (0,0) {};
\node[ptsq={myblue}{}] at (\xL,{\f(\xL)}) {};

\foreach \ta in {0.1,0.2,...,0.9}
{
\pgfmathsetmacro{\ya}{\f(\ta)}
\pgfmathsetmacro{\slopea}{\fp(\ta)}
\pgfmathsetmacro{\txa}{1/sqrt(1+\slopea^2)}
\pgfmathsetmacro{\tya}{\slopea/sqrt(1+\slopea^2)}
\pgfmathsetmacro{\nxa}{-\slopea/sqrt(1+\slopea^2)}
\pgfmathsetmacro{\nya}{1/sqrt(1+\slopea^2)}
\pgfmathsetmacro{\xx}{\ta-\ya*\nxa/\nya}

\coordinate (P) at (\ta,\ya);

\draw[semithick] (P) -- (\xx, 0);

\def\eps{0.01}

\draw[semithick]
  (P)
  --
  ($(P)!\eps cm!({\ta+\txa},{\ya+\tya})$)
  --
  ($($(P)!\eps cm!({\ta+\txa},{\ya+\tya})$)!\eps cm!90:(P)$)
  --
  ($(P)!\eps cm!-90:({\ta+\txa},{\ya+\tya})$)
  --
  cycle;
}


\def\ta{0.3}

\pgfmathsetmacro{\ya}{\f(\ta)}
\pgfmathsetmacro{\slopea}{\fp(\ta)}
\pgfmathsetmacro{\txa}{1/sqrt(1+\slopea^2)}
\pgfmathsetmacro{\tya}{\slopea/sqrt(1+\slopea^2)}
\pgfmathsetmacro{\nxa}{-\slopea/sqrt(1+\slopea^2)}
\pgfmathsetmacro{\nya}{1/sqrt(1+\slopea^2)}
\pgfmathsetmacro{\xx}{\ta-\ya*\nxa/\nya}

\coordinate (P) at (\ta,\ya);

\draw[myblue] plot[domain=0:\xL,samples=20,smooth] ({\x},{\f(\x)});

\draw[myred] (P) -- (\xx, 0);
\node[ptcl={myred}{above:{$\projvariete(x)$}}] at (P) {};
\node[ptcl={myred}{below:{$x$}}] at (\xx, 0) {};


\end{tikzpicture}
    \caption{Projection orthogonale sur la surface lisse $\surflisse$.}
    \label{fig:estimation_projection}
\end{figure}

\begin{lemme}[{\cite[Lemme 4.1]{dziuk_finite_2013}}] \label{lemme:projection_dziuk}
Supposons que $\surflisse$ et $\surfpoly$ soient telles que décrites dans la \Cref{sec:triangulations}. L'application $\fonctionens[\projvariete]{\surfpoly}{\surflisse}$ est bijective (voir le \Cref{lemme:fermi_coordinates}). Pour la fonction distance orientée à $\surflisse$, définie par la \Cref{eq:def_oriented_distance}, nous avons l'estimation
\[
\norme{d}_{\L^\infty(\surfpoly)} \leqslant c \size(\ensprimalplat)^2.
\]
Afin d'alléger les notations, nous posons $h = \size(\ensprimalplat)$. Le quotient $\delta_h$ entre les mesures surfaciques lisse et discrète $\d A$ et $\d A_h$, défini par $\delta_h \d A_h = \d A$, vérifie
\[
\norme{1 - \delta_h}_{\L^\infty(\surfpoly)} \leqslant c \size(\ensprimalplat)^2.
\]
\end{lemme}

\begin{remarque} \label{remarque:sizeM_vs_sizeT}
D'après la définition \ref{eq:size_ensddfvplat} de $\size(\ensddfvplat)$, on a $\size(\ensprimalplat) \leqslant \size(\ensddfvplat)$ donc on peut écrire les majorations du \Cref{lemme:projection_dziuk} en fonction de $\size(\ensddfvplat)$.  
\end{remarque}

\begin{demo}[\Cref{prop:estimations_tailles_plat_vs_courbe}]
On suppose $\size(\ensddfvplat)$ suffisamment petit de sorte que la surface polyédrique $\surfpoly$ soit dans le voisinage tubulaire $U_\delta$ défini dans le \Cref{lemme:fermi_coordinates}. On peut alors appliquer le \Cref{lemme:projection_dziuk} et la \Cref{remarque:sizeM_vs_sizeT} d'après lesquels le quotient $\delta_h$ entre les mesures surfaciques lisse et polyédrique vérifie $\norme{1 - \delta_h}_{\L^\infty(\surfpoly)} \leqslant c \size(\ensddfvplat)^2$, où $c > 0$ est indépendante de $\size(\ensddfvplat)$. Ainsi, $\delta_h \geqslant 1 - c \size(\ensddfvplat)^2$ presque partout sur $\surfpoly$. Soit $E \subset \surfpoly$ une partie mesurable. Puisque la projection $\fonctionens[\projvariete]{\surfpoly}{\surflisse}$ est bijective, la formule de changement de variables donne
\[
\mesure(\projvariete E)
=
\int_E \delta_h \d A_h
\geqslant
\big(1 - c \size(\ensddfvplat)^2\big) \mesure(E).
\]
En appliquant cette inégalité successivement à $E = \primalplat{K}$ et à $E = \mailledualeplate{K}$, et en utilisant $\primal{K} = \projvariete \primalplat{K}$ et $\mailleduale{K} = \projvariete \mailledualeplate{K}$, on obtient le résultat.
\end{demo}

\begin{demo}[{\Cref{prop:mesure_sigmacourbe}}]
On suppose $\size(\ensddfvplat)$ suffisamment petit de sorte que la surface polyédrique $\surfpoly$ soit dans le voisinage tubulaire $U_\delta$ défini dans le \Cref{lemme:fermi_coordinates}. Considérons un segment $\ell \subset \surfpoly$ paramétré par $\fonctionens[\gamma]{\intervalleff{0}{\mesure(\ell)}}{\ell}$. Le segment courbe associé $\projvariete \ell$ est quant à lui paramétré par $\gamma_\surflisse = \projvariete{} \circ \gamma$  (voir \Cref{fig:estimation_projection}). L'objectif est d'estimer la longueur $\mesure(\projvariete \ell) = \int_0^{\mesure(\ell)} \abs{\gamma'_\surflisse(s)} \d s$ en fonction de $\mesure(\ell)$. Or $\abs{\gamma'_\surflisse(s)} = \abs{\diff \projvariete\mathopen{}\big(\gamma(s)\big) \, \gamma'(s)}$. Nous allons donc chercher à estimer la différentielle $\diff \projvariete$. D'après le \Cref{lemme:fermi_coordinates}, on a $\projvariete(x) = x - d(x) \grad d(x)$. En différentiant dans la direction $v \in \R^3$, on a pour $x \in U_\delta$,
\begin{equation} \label{eq:diff_projvariete}
\diff \projvariete(x) \, v = v - (\grad d(x) \cdot v ) \grad d(x) - d(x) \, \diff^2 d(x) \, v = \big( P(x) - d(x) \, H(x) \big) v
\end{equation}
avec
\[
P(x) \defeq \I - \grad d(x) \otimes \grad d(x) 
\quad
\text{et}
\quad
H(x) \defeq \diff^2 d(x).
\]
La matrice $P(x)$ est la projection orthogonale sur l'espace tangent à la surface de niveau de $d$ passant par $x$. Ce plan est parallèle à $T_{\projvariete(x)} \surflisse$. On déduit de \ref{eq:diff_projvariete} que $\diff \projvariete(x) = P(x) - d(x) \, H(x)$ et comme $P(x)$ est une projection orthogonale, $\norme{P(x)}=1$ ce qui entraîne
\begin{equation} \label{eq:estim_diff_proj}
\norme{\diff \projvariete(x)} \leqslant 1 + \abs{d(x)} \norme{H(x)} \quad \forall x \in U_\delta.
\end{equation}
D'après le \Cref{lemme:fermi_coordinates}, on a $d \in \Cont^2(U_\delta)$ donc l'application $x \mapsto H(x)$ est continue sur $U_\delta$. Soit $\delta' \in \intervalleoo{0}{\delta}$. On suppose $\size(\ensddfvplat)$ suffisamment petit pour que $\surfpoly \subset U_{\delta'}$. Alors $\overline{U_{\delta'}} \subset U_\delta$ est compact et $H$ y est continu. On peut alors poser $C_H \defeq \norme{H}_{\L^\infty(\overline{U_{\delta'}})} < + \infty$. Nous appliquons la première estimation du \Cref{lemme:projection_dziuk} et la \Cref{remarque:sizeM_vs_sizeT} à \ref{eq:estim_diff_proj} pour obtenir, pour tout $x \in \surfpoly \subset U_\delta$,
\[
\norme{\diff \projvariete(x)} \leqslant 1 + C \size(\ensddfvplat)^2,
\quad   
C \defeq c \, C_{H}.
\]
La constante $C$ dépend de la géométrie de $\surflisse$ mais pas de la taille du maillage. Nous pouvons maintenant revenir au calcul de $\mesure(\projvariete \ell)$ : 
\[
\mesure(\projvariete \ell)
= \int_0^{\mesure(\ell)} \abs{\diff \projvariete\mathopen{}\big(\gamma(s)\big) \, \gamma'(s)} \d s \leqslant \big( 1 + C \size(\ensddfvplat)^2 \big) \mesure(\ell).
\]
\end{demo}